\chapter{Závěr}
\label{sec:conclusion}

Cílem této bakalářské práce bylo navrhnout a implementovat systém
chování nehráčských postav (NPC) v herním prostředí s využitím
behaviorálních stromů v prostředí Unreal Engine~5.

Navržený systém kombinuje využití \texttt{AI Controlleru},
\texttt{Blackboardu} a behaviorálního stromu, což umožňuje
modulární a přehlednou implementaci jednotlivých typů chování.
Součástí řešení je také systém percepce založený na zrakových
a sluchových podnětech a mechanismus sdílení informací o
předmětech prostřednictvím \texttt{GameState}.

V rámci implementace se podařilo vytvořit funkční prototyp,
ve kterém NPC dokáží hlídkovat, reagovat na hráče,
interagovat s objekty ve scéně a pracovat s předměty.
Chování jednotlivých postav odpovídá návrhu a je snadno
rozšiřitelné o další prvky.

Testování ukázalo, že základní mechanismy systému fungují
spolehlivě, zejména v oblasti percepce, přechodů mezi stavy
a pohybu po hlídkových trasách. Naopak některé pokročilejší
části, jako manipulace s předměty nebo sdílení informací mezi
NPC, nebyly plně ověřeny a představují prostor pro další
vylepšení.

\section{Možná rozšíření a budoucí práce}

Navržený systém chování NPC by bylo možné dále rozšířit o reakce na dynamické změny prostředí. 
Jednou z možností je zohlednění denní doby a povětrnostních podmínek.
Například v nočních hodinách by NPC mohly trávit více času v interiéru, zatímco během dne by se častěji pohybovaly v okolí domu. 
Podobně by bylo možné simulovat vliv deště, který by negativně ovlivňoval percepční schopnosti NPC, například zhoršením sluchu.

Dalším rozšířením by bylo zavedení interaktivních prvků v prostředí, jako jsou například světelné vypínače.
NPC by mohly reagovat na změny osvětlení v místnostech, například snahou rozsvítit při vstupu do tmavého prostoru, nebo naopak vyhýbáním se neosvětleným místnostem.
Takové chování by mohlo přispět k vyšší věrohodnosti a variabilitě simulace.

Uvedená rozšíření nebyla v rámci této práce implementována, je ale možné takových cílů dosáhnout při dalším vývoji této hry.

\endinput